% why/motivation.tex

\chapter{Motivation}

\section{Why a Railway Transition Is Not a Road-Corridor Detail}

A road corridor influences a driver's available path, but a railway vehicle is
kinematically constrained by its rail pair. The track therefore prescribes the
curvature history experienced by the vehicle. A change in
\(\kappa(s)\), its derivatives, cant \(u(s)\), or speed is transmitted into
wheel--rail interaction, vehicle response, passenger comfort, track loading,
wear, and safety-related assessment. Transition design is consequently a
central railway problem, not merely a convenient interpolation between map
elements.

Consider two laws with the same length and boundary curvatures. A clothoid has
constant curvature derivative. A polynomial or trigonometric half-wave can
instead arrange selected endpoint conditions on
\(\kappa'(s)\) and \(\kappa''(s)\). Similar-looking centerlines can therefore
carry different dynamic implications. Conversely, demanding smoother endpoint
behaviour may require more length or a different allocation of degrees of
freedom. The clothoid is important, simple, and frequently compact; these
advantages do not make it the universal physical optimum
\cite{BrustadDalmo2020Review,Koc2017SmoothedEnds}.

This comparison produced the originating research direction of the work.
Berlinish sought a shared functional grammar for such alternatives, while
axtranNew sought to calculate their admissible compositions. AIM later widened
the question from \emph{which transition can be calculated?} to \emph{how can
the resulting engineering knowledge retain meaning and become usable?}

\section{The Curvature Edit}

The curvature edit introduced in the opening moves the downstream positions and
changes the tangent direction. If cant and speed remain fixed, vehicle-related
quantities change as well. Constraints satisfied before the edit may no longer
be satisfied afterwards; standard railway design couples curvature, cant,
speed, and their rates of change \cite{EN13803_2017,Klein1998Trassierung}.

The visible curve is therefore the beginning of the consequence, not its end.
The edit must be interpreted against terrain and adjacent assets, compared with
the existing or measured track, checked against applicable rules, and presented
as an alternative for review. At every step, a different representation may be
appropriate. What must not disappear is the relation among the alignment, the
edit, the resulting states, the evidence, and the authority to decide.

\section{Three Powerful Views, Three Necessary Limits}

GIS, CAD, and BIM approach this situation from different strengths. GIS places
the alignment in extensive spatial context and supports interoperable access,
processing, visualization, and discovery of geospatial information
\cite{OGCStandardsProgram2025}. This breadth is indispensable for land,
environment, networks, and spatial analysis, but the engineer must still
determine which available context is applicable to the alignment change.

CAD concentrates geometric construction, communication, and direct
manipulation \cite{KimKumarSnider2015Geometry}. It makes the curvature edit
tangible and produces geometry needed for design and delivery. Its geometric
model, however, does not by itself establish the provenance of an observation,
the authority behind a rule, or the durable identity shared by another
representation.

BIM organizes built-asset information for exchange among software applications
and project participants. IFC makes this role explicit through schemas and
workflow-specific model views \cite{BuildingSMARTIFC4x3Scope}. These structures
can carry alignment information without thereby owning its intrinsic
construction or all knowledge required to evaluate a change.

The issue is not to choose a winner among the three. The issue is that the
alignment must pass between their views without becoming merely a GIS feature,
a CAD curve, or a BIM container. A representation is useful precisely because
it selects and expresses something for a purpose. That selection also means no
single representation should silently become the whole engineering object.

\section{Fragmented Knowledge at the Moment of Action}

Some relevant knowledge is encoded in software. Some resides in standards,
project requirements, measurement reports, calculation procedures, and design
decisions. Some remains with engineers who know which rule is applicable and
which apparent discrepancy is consequential. Fragmentation becomes costly when
an action demands these pieces together. Construction-informatics research
relates this persistent fragmentation to specialization, heterogeneous tools,
and only partially interoperable collaboration \cite{Turk2020Interoperability}.

For the curvature edit, the engineer needs more than data access. The software
must know which object is being changed, which state is intended, which
observations describe another state, which context selects the rules, and which
quantities depend on the edit. It must also preserve the boundary between
calculation and authorization.

This is the practical motivation for making engineering knowledge applicable.
Applicable knowledge is not simply text stored beside geometry. It has scope,
provenance, conditions, and an accountable meaning; it must be connected to the
state and operation for which it matters.

\section{Why Curvature Exposes the Problem}

A railway alignment is often introduced as a sequence of straights, circular
arcs, and transition curves. That description is useful, but curvature reveals
the generative relation beneath it. Along intrinsic arc length (s), curvature
governs the change of heading,
\[
  \theta'(s)=\kappa(s),
\]
and integration produces planar pose and position. Changes in curvature and
cant also enter engineering assessments of vehicle response and comfort
\cite{Kufver2000RailwayAlignment,BrustadDalmo2020Review}.

Curvature thus connects constructive intent with geometric and dynamic
consequence. Yet the resulting coordinates still do not own the alignment's
identity, and a measured curve still does not automatically reveal the intended
curvature law. The distinction between intrinsic construction, metric
realization, physical realization, and observation is not philosophical
decoration. It determines what can safely be inferred after the edit.

\section{The Requirement for AIM}

A useful response must satisfy four demands at once:

\begin{itemize}
\item preserve alignment identity across purposeful representations;
\item expose the constructive and state dependencies through which changes
  acquire consequences;
\item keep intended, realized, and observed information distinct and traceable;
\item apply engineering knowledge without assigning decision authority to a
  file format, software component, or numerical solver.
\end{itemize}

Alignment-Based Information Modelling is developed to meet these demands. Its
starting point is not a larger file format. It is an alignment-specific account
of what the engineering object is, how it acquires measurable and operational
state, how evidence relates to it, and how knowledge participates in
evaluation.

\begin{summary}
The difficulty exposed by a curvature edit is not that software cannot draw the
new curve. It is that the engineering meaning of the change crosses tools,
states, evidence, rules, and authority. AIM begins by preserving those
distinctions and making their dependencies explicit.
\end{summary}
